% ==========================================================
%  EXAM QUESTIONS — AI IN LIFE SCIENCES
%  Article class, shared visual language with summary.tex.
% ==========================================================

\documentclass[11pt,oneside]{article}

\usepackage[
  a4paper,
  top=2.4cm,
  bottom=2.1cm,
  left=2.3cm,
  right=2.3cm,
  headheight=15pt
]{geometry}

\usepackage[T1]{fontenc}
\usepackage[utf8]{inputenc}
\usepackage{lmodern}
\usepackage{microtype}
\usepackage{inconsolata}
\usepackage{amsmath,amssymb,mathtools,bm}
\usepackage{enumitem}
\usepackage{booktabs,array,multicol,longtable}
\usepackage{xcolor}
\usepackage[most]{tcolorbox}
\usepackage{titlesec}
\usepackage{fancyhdr}
\usepackage[hidelinks]{hyperref}

\setlength{\parindent}{0pt}
\setlength{\parskip}{0.55em}
\setlist[itemize]{topsep=3pt,itemsep=2pt,leftmargin=2.2em}
\setlist[enumerate]{topsep=3pt,itemsep=2pt,leftmargin=2.2em}

% ── Colors ───────────────────────────────────────────────
% Shared accent with summary.tex — do not change.
\definecolor{accent}{HTML}{1B3A5C}
\definecolor{q@frame}{HTML}{1A5276}
\definecolor{q@bg}{HTML}{F6FAFD}
\definecolor{key@frame}{HTML}{1E8449}
\definecolor{key@bg}{HTML}{F4FBF6}
\definecolor{note@frame}{HTML}{9A6A00}
\definecolor{note@bg}{HTML}{FEF9E7}

% ── Box global style ─────────────────────────────────────
\tcbset{
  enhanced,
  breakable,
  arc=2.5mm,
  boxrule=0.5pt,
  left=4.5mm,
  right=4mm,
  top=2.5mm,
  bottom=2.5mm,
  fonttitle=\bfseries\small,
  sharp corners=west,
  parbox=false
}

% Question box (blue) — one per exam question.
\newtcolorbox{qbox}[1][]{
  colback=q@bg,
  colframe=q@frame,
  borderline west={3.5pt}{0pt}{q@frame},
  title={#1}
}

% Key insight box (green) — answer hints, key results.
\newtcolorbox{keybox}[1][]{
  colback=key@bg,
  colframe=key@frame,
  borderline west={3.5pt}{0pt}{key@frame},
  title={#1}
}

% Organizational note box (amber) — structural remarks only.
\newtcolorbox{notebox}[1][]{
  colback=note@bg,
  colframe=note@frame,
  borderline west={3.5pt}{0pt}{note@frame},
  title={#1}
}

% ── Section headings ─────────────────────────────────────
\titleformat{\section}
  {\normalfont\Large\bfseries\color{accent}}
  {}{0em}{}
  [{\vspace{1pt}\color{accent!45}\titlerule[0.45pt]}]
\titlespacing*{\section}{0pt}{3.2ex plus 1ex minus .2ex}{1.8ex plus .2ex}

\titleformat{\subsection}
  {\normalfont\normalsize\bfseries\color{accent!80!black}}
  {}{0em}{}
\titlespacing*{\subsection}{0pt}{2.2ex plus 1ex minus .2ex}{1.0ex plus .2ex}

% ── Header & footer ──────────────────────────────────────
\pagestyle{fancy}
\fancyhf{}
\fancyhead[L]{\small\color{gray!70}Old Exam Questions — \coursetitle}
\fancyhead[R]{\small\color{gray!70}\thepage}
\renewcommand{\headrulewidth}{0.4pt}

% ── Document metadata ─────────────────────────────────────
\newcommand{\coursetitle}{AI in Life Sciences}

% ── Question counter & macros ────────────────────────────
\newcounter{question}
\newcommand{\questiontitle}[1]{%
  \stepcounter{question}\begin{qbox}[Question \thequestion: #1]}
\newcommand{\closequestion}{\end{qbox}}

% Mark questions that appeared on multiple exams.
\newcommand{\repeatnote}{\textit{\color{gray!70}Asked multiple times.}}

% Multiple-choice checkbox bullet.
\newcommand{\choice}{\item[$\square$] }

% Fill-in answer line. Usage: \answerline{3cm}
\newcommand{\answerline}[1]{\makebox[#1]{\hrulefill}}

% Blank working area — dashed-border box for student calculations.
\newcommand{\workingspace}[1]{%
  \par\vspace{2pt}%
  \begin{tcolorbox}[
    enhanced,
    breakable=false,
    colback=white,
    colframe=gray!40,
    boxrule=0.45pt,
    arc=0pt,
    sharp corners,
    height=#1,
    left=3mm, right=3mm, top=3mm, bottom=3mm,
    frame style={dashed}
  ]%
  \end{tcolorbox}%
  \vspace{2pt}%
}

% ── Math macros (minimal subset for exam questions) ──────
\newcommand{\R}{\mathbb{R}}
\newcommand{\E}{\mathbb{E}}
\newcommand{\Prob}{\mathbb{P}}
\newcommand{\KL}[2]{D_{\mathrm{KL}}\!\left(#1\,\|\,#2\right)}
\newcommand{\dd}{\,\mathrm{d}}


% ==========================================================
%  BEGIN DOCUMENT
% ==========================================================
\begin{document}

% ── Title page ───────────────────────────────────────────
\thispagestyle{empty}
\vspace*{1.4cm}
\begin{center}
  {\color{accent}\rule{\linewidth}{1.8pt}}\par
  \vspace{0.65cm}
  {\fontsize{26}{32}\selectfont\bfseries\color{accent}Old Exam Questions\\[2pt] \coursetitle\par}
  \vspace{0.55cm}
  {\color{accent!35}\rule{\linewidth}{0.5pt}}\par
  \vspace{0.55cm}
  {\large Stefan Eder\par}
  \vspace{0.15cm}
  {\small\color{gray!70}\today\par}
\end{center}

\vspace{0.9cm}
\begin{notebox}[How this file is organized]
Questions come from two sittings: the \textbf{2024} exam (25 June 2024,
Sections~1--6) and the \textbf{2026} exam (25 June 2026, Sections~7--10). The 2024
items are \emph{multiple-select} (``which of the following is/are true'' means
\textbf{one or more} of the four options are correct). The 2026 items are
\emph{single-best-answer} (exactly \textbf{one} of the three options is correct).
To avoid reproducing the originals verbatim, both the question wording and the
answer options have been paraphrased (the meaning and the correct answers are
preserved). Questions are grouped by topic rather than in their original random
order. The full answer key is at the
end --- the question boxes themselves are intentionally left blank so the file
doubles as a practice set.
\end{notebox}

\tableofcontents
\newpage


% ==========================================================
%  SECTION 1 — Drug Discovery, Molecules & Representations
% ==========================================================
\section{Drug Discovery, Molecules \& Representations}

\questiontitle{SMILES strings}
Regarding the SMILES notation for molecules, which of the following statements hold?
\begin{itemize}[leftmargin=1.8em]
  \choice SMILES is short for the \emph{Simplified Molecular Input Line Entry System}.
  \choice A SMILES encodes a molecule as a one-dimensional sequence of characters.
  \choice A SMILES string cannot be turned into a molecular graph.
  \choice Each molecule has exactly one SMILES, so no two distinct molecules can share a string.
\end{itemize}
\closequestion

\questiontitle{Reading a SMILES string}
Consider the molecule \emph{pyruvic acid} (2-oxopropanoic acid), i.e.\ a methyl
group followed by a ketone carbonyl and then a carboxylic acid
($\text{CH}_3\text{--C}(\text{=O})\text{--C}(\text{=O})\text{--OH}$). Which of the
SMILES strings below correctly encode this molecule?
\begin{itemize}[leftmargin=1.8em]
  \choice \texttt{O=C(C(=O)O)C}
  \choice \texttt{CC(=O)C(=O)O}
  \choice \texttt{CCCOOO}
  \choice \texttt{OCCOO}
\end{itemize}
\closequestion

\questiontitle{Molecular fingerprints}
Which of the following statements about molecular fingerprints are accurate?
\begin{itemize}[leftmargin=1.8em]
  \choice A unique molecular graph can be reconstructed from them.
  \choice They are binary or numeric encodings of molecular structure.
  \choice They serve similarity searches and machine-learning tasks.
  \choice They give a 3D representation of the molecule.
\end{itemize}
\closequestion

\questiontitle{Bioassays}
Select the correct statements about bioassays.
\begin{itemize}[leftmargin=1.8em]
  \choice They quantify the biological response or function that a substance triggers.
  \choice They usually probe the substance across a range of concentrations.
  \choice They are wet-lab procedures for characterising a compound in a biological setting.
  \choice They are usually carried out \emph{in vitro}, i.e.\ inside a living organism.
\end{itemize}
\closequestion

\questiontitle{ADME properties}
Which statements about ADME (Absorption, Distribution, Metabolism, and Excretion)
properties are correct?
\begin{itemize}[leftmargin=1.8em]
  \choice Lipinski's rule of five can serve as a proxy for ADME behaviour.
  \choice Favourable ADME properties guarantee that a compound becomes an effective drug.
  \choice Certain ADME properties, such as clearance, are hard to model.
  \choice ADME properties are decisive for a drug's pharmacokinetics.
\end{itemize}
\closequestion

\questiontitle{Characteristics of QSAR / bioactivity data}
Which of the following correctly characterise typical QSAR / bioactivity datasets?
\begin{itemize}[leftmargin=1.8em]
  \choice QSAR datasets typically contain on the order of thousands to millions of data points.
  \choice QSAR datasets usually show a balanced mix of active and inactive compounds.
  \choice QSAR datasets frequently include groups of structurally related compounds.
  \choice QSAR labels are often noisy because of biological variability and experimental inconsistency.
\end{itemize}
\closequestion

\questiontitle{Applications of QSAR modeling}
For which of the following applications is QSAR modeling relevant?
\begin{itemize}[leftmargin=1.8em]
  \choice Building 3D models of protein structures.
  \choice Finding promising drug candidates.
  \choice Estimating the toxicity of novel chemical compounds.
  \choice Mapping the genome sequences of organisms.
\end{itemize}
\closequestion


% ==========================================================
%  SECTION 2 — Virtual Screening & QSAR Modeling
% ==========================================================
\section{Virtual Screening \& QSAR Modeling}

\questiontitle{Ligand-based virtual screening}
Which statements about ligand-based virtual screening are correct?
\begin{itemize}[leftmargin=1.8em]
  \choice Molecular fingerprints are a common ingredient of ligand-based virtual screening.
  \choice Ligand-based virtual screening cannot predict the bioactivity of untested compounds.
  \choice Ligand-based virtual screening needs the 3D structure of the target protein.
  \choice Ligand-based virtual screening selects compounds using bioactivity predictions alone.
\end{itemize}
\closequestion

\questiontitle{Data splits in QSAR model evaluation}
When evaluating QSAR models, which of the following statements about
data-splitting strategies hold?
\begin{itemize}[leftmargin=1.8em]
  \choice Temporal splits gauge performance on future, unseen data by mimicking how data accrue over time.
  \choice A random split is generally the best evaluation scheme because it removes all bias.
  \choice Cross-validation is seldom applied to QSAR because of its computational cost.
  \choice Cluster splits keep similar compounds together so that structurally similar ones do not span the training and test sets.
\end{itemize}
\closequestion

\questiontitle{Early-recognition metrics in virtual screening}
Because only the top of a virtual-screening ranking is followed up experimentally,
recognising actives early matters most. Which of the following metrics are
specifically designed to reward strong early recognition?
\begin{itemize}[leftmargin=1.8em]
  \choice Precision-Recall Curve.
  \choice BEDROC (Boltzmann-Enhanced Discrimination of ROC).
  \choice Average Precision (AP).
  \choice Enrichment Factor (EF).
\end{itemize}
\closequestion

\questiontitle{Target identification / deconvolution}
Which statements about target identification / deconvolution are correct?
\begin{itemize}[leftmargin=1.8em]
  \choice For a given molecule, it determines which of several simultaneously present proteins it is most likely to bind.
  \choice It can be approached with multi-class classification.
  \choice It can be approached with multi-task classification.
  \choice For a given molecule, it determines every protein it would bind when several are present at once.
\end{itemize}
\closequestion

\questiontitle{Drug synergy and antagonism}
Which of the following statements about drug synergy and antagonism are correct?
\begin{itemize}[leftmargin=1.8em]
  \choice Synergy means the combined effect of two or more drugs exceeds the sum of their individual effects.
  \choice The Loewe additivity model is a standard tool for assessing synergy.
  \choice Antagonism means the combined effect of two or more drugs exceeds the sum of their individual effects.
  \choice Synergy/antagonism analysis applies only to chemically similar drugs.
\end{itemize}
\closequestion

\questiontitle{Random Forests vs.\ Graph Neural Networks}
Comparing Random Forests (RF) and Graph Neural Networks (GNNs) for molecular
modeling, which statements are correct?
\begin{itemize}[leftmargin=1.8em]
  \choice GNNs are inherently superior for QSAR because they capture the graph structure of molecules, giving more accurate predictions in every case.
  \choice RF models overfit less than GNNs, which makes them more reliable on datasets with many features but few samples.
  \choice RF models are usually more interpretable than GNNs, so the contribution of individual molecular features is easier to read off.
  \choice GNNs are less sensitive to input-data quality and variability than RFs, making them more consistent across different datasets.
\end{itemize}
\closequestion


% ==========================================================
%  SECTION 3 — Graph Neural Networks & Advanced Methods
% ==========================================================
\section{Graph Neural Networks \& Advanced Methods}

\questiontitle{Molecular graphs and graph neural networks}
Which statements about molecular graphs and graph neural networks are correct?
\begin{itemize}[leftmargin=1.8em]
  \choice A molecular graph can be stored as two matrices: one for the node features and one for the edges between nodes.
  \choice Node features may include the atomic number, hybridization state, and formal charge.
  \choice A GNN uses only the node features to form its final prediction.
  \choice Bond-type information (e.g.\ single vs.\ double bond) cannot be represented in a GNN.
\end{itemize}
\closequestion

\questiontitle{Limitations of GNNs in molecular modeling}
Which of the following are genuine limitations of GNNs in the context of
molecular modeling?
\begin{itemize}[leftmargin=1.8em]
  \choice GNNs cannot predict any molecular property with accuracy.
  \choice GNNs may struggle to capture long-range interactions between atoms.
  \choice GNNs can have trouble with very large, complex molecular graphs.
  \choice GNNs need considerable computational resources to train.
\end{itemize}
\closequestion

\questiontitle{Multi-modal contrastive learning}
Which statements about multi-modal contrastive learning are correct?
\begin{itemize}[leftmargin=1.8em]
  \choice Multi-modal contrastive learning embeds several modalities (e.g.\ molecules and proteins) into one shared latent space.
  \choice Multi-modal contrastive learning is a self-supervised method.
  \choice Multi-modal contrastive learning concentrates the (latent) representations of several modalities (e.g.\ molecules and proteins) to make predictions.
  \choice Multi-modal contrastive learning is an unsupervised method.
\end{itemize}
\closequestion

\questiontitle{Multiple instance learning (MIL)}
Which of the following statements about multiple instance learning (MIL) hold?
\begin{itemize}[leftmargin=1.8em]
  \choice MIL trains on sets (bags) of instances in which each instance carries its own label.
  \choice MIL needs the bag size to be fixed.
  \choice MIL trains on sets (bags) of instances that share a single label.
  \choice A MIL model should be invariant to the ordering of the instances.
\end{itemize}
\closequestion


% ==========================================================
%  SECTION 4 — Generative Models & Chemical Synthesis
% ==========================================================
\section{Generative Models \& Chemical Synthesis}

\questiontitle{Generative models for small molecules}
Which statements about generative models for small molecules are correct?
\begin{itemize}[leftmargin=1.8em]
  \choice Because SMILES syntax is intricate with long-range dependencies, SMILES generators yield fewer than 10\% valid strings.
  \choice Generative models are used to explore the already-synthesised, known chemical space.
  \choice Generative adversarial networks cannot generate small molecules.
  \choice Autoregressive models can generate SMILES representations of small molecules.
\end{itemize}
\closequestion

\questiontitle{Goal-directed molecule generation}
Which of the following correctly describe goal-directed molecule generation?
\begin{itemize}[leftmargin=1.8em]
  \choice Goal-directed generation is concerned only with reproducing existing structures that already have the desired properties.
  \choice Goal-directed generation aims to produce molecules whose properties resemble those of the training set.
  \choice Reinforcement learning can be used to drive goal-directed molecule generation.
  \choice Goal-directed generation aims to produce molecules with particular, pre-specified properties.
\end{itemize}
\closequestion

\questiontitle{Fr\'echet ChemNet Distance (FCD)}
Which statements about the Fr\'echet ChemNet Distance (FCD) and how it is computed
are correct?
\begin{itemize}[leftmargin=1.8em]
  \choice FCD measures the distance between the individual atoms of molecules.
  \choice FCD compares the means and covariances of the molecules' latent representations.
  \choice FCD is computed straight from the raw molecular structures with no transformation.
  \choice FCD is a metric comparing the distribution of generated molecules with that of real ones.
\end{itemize}
\closequestion

\questiontitle{AI for chemical reactions}
Which of the following statements about machine learning for chemical reactions are
correct?
\begin{itemize}[leftmargin=1.8em]
  \choice Some approaches use so-called reaction templates, i.e.\ transformation rules for molecules.
  \choice Deep-learning reaction-outcome predictors can be trained on databases of known reactions.
  \choice Transformer architectures have been applied to predicting reaction outcomes.
  \choice Predicting reactions with machine learning requires quantum-mechanical data because of bond changes and electron movement.
\end{itemize}
\closequestion

\questiontitle{Synthesis planning}
Which statements about synthesis planning are correct?
\begin{itemize}[leftmargin=1.8em]
  \choice Synthesis planning designs a sequence of chemical reactions that yields a target compound.
  \choice Retrosynthetic analysis is central to it, decomposing complex molecules into simpler starting materials.
  \choice Modern deep-learning methods have made synthesis planning trivial.
  \choice Synthesis planning amounts to nothing more than choosing the reactants.
\end{itemize}
\closequestion

\questiontitle{Retrosynthesis as a reinforcement-learning game}
Multi-step retrosynthesis can be cast as a game solvable by reinforcement learning.
Which of the following statements about this ``retrosynthesis game'' are correct?
\begin{itemize}[leftmargin=1.8em]
  \choice In the ``game'', an action consists of choosing reactants for a reaction.
  \choice In the ``game'', the state is the set of available building blocks.
  \choice In the ``game'', the environment carries out the reaction and returns the new molecules.
  \choice In the ``game'', the reward is granted once the product has been synthesised from the building blocks.
\end{itemize}
\closequestion


% ==========================================================
%  SECTION 5 — Medical Data & Imaging
% ==========================================================
\section{Medical Data \& Imaging}

\questiontitle{Tabular medical data: types and missing values}
Suppose we are given a table of patient medical records with three columns:
``smoking status'' (one of three string values: non-smoker, former smoker,
current smoker), ``cholesterol level'' (floating-point measurements in mg/dL), and
``diabetes risk score'' (integer values in $[0,10]$ indicating the likelihood of
developing diabetes). Which of the following statements are correct?
\begin{itemize}[leftmargin=1.8em]
  \choice Missing ``cholesterol level'' entries can be filled with the column's median to cope with skew.
  \choice The ``smoking status'' column is continuous, since it reflects patients' various smoking habits.
  \choice ``cholesterol level'' is ordinal, because the measurements can be placed in increasing order.
  \choice Missing ``smoking status'' entries can be filled with the column's mean to cope with skew.
\end{itemize}
\closequestion

\questiontitle{Batch effects in medical data}
Which statements about batch effects in medical data are correct?
\begin{itemize}[leftmargin=1.8em]
  \choice Overlooking batch effects can yield models that separate batches instead of the true biological conditions.
  \choice Batch effects vanish simply by collecting more samples.
  \choice Batch effects introduce artificial differences that stem from varying experimental conditions.
  \choice Batch effects are generally helpful, supplying needed variability to the data.
\end{itemize}
\closequestion

\questiontitle{Domain shifts in medical and healthcare data}
Medical and healthcare data are frequently subject to domain shifts. Using the
lecture notation ($y$ = disease status, e.g.\ ``healthy''/``diseased'';
$x$ = patient features, e.g.\ age, blood parameters, heart rate), which of the
following statements about domain shifts are correct?
\begin{itemize}[leftmargin=1.8em]
  \choice A model with high predictive quality on a randomly drawn validation set will, by that fact alone, adapt to domain shifts and keep its quality over time.
  \choice The patient-feature distribution $p(x)$ is never touched by domain shifts.
  \choice During waves of an infectious disease (such as COVID-19), the prevalence---and thus the probability $p(y)$ of observing the phenotype---changes.
  \choice If the way the status $y$ is diagnosed changes, e.g.\ an antigen test is replaced by a PCR test, this is a concept shift $p(y\mid x)$.
\end{itemize}
\closequestion

\questiontitle{Medical imaging modalities}
Which statements about medical imaging modalities are correct?
\begin{itemize}[leftmargin=1.8em]
  \choice Computed Tomography (CT) builds its body images from ultrasound.
  \choice X-rays distinguish tissues of differing density, such as bone versus soft tissue.
  \choice Magnetic Resonance Imaging (MRI) yields detailed images of internal structures using strong magnetic fields and radio waves.
  \choice Optical coherence tomography forms images of organs and tissues using multiple X-rays.
\end{itemize}
\closequestion

\questiontitle{Most common task in medical imaging}
Compared with typical computer-vision tasks, which task is especially prominent in
medical imaging?
\begin{itemize}[leftmargin=1.8em]
  \choice Image generation.
  \choice Classification.
  \choice Segmentation.
  \choice Object detection.
\end{itemize}
\closequestion

\questiontitle{The U-Net architecture}
Which of the following correctly describe the U-Net architecture?
\begin{itemize}[leftmargin=1.8em]
  \choice U-Net was the first vision model whose encoder relied on self-attention rather than convolutions.
  \choice The U-Net architecture has a symmetric encoder and decoder.
  \choice U-Net was created specifically for the object-detection task.
  \choice To help decoding, encoder feature maps are concatenated with those of the matching up-sampling level of the decoder.
\end{itemize}
\closequestion

\questiontitle{The ``Clever Hans'' effect}
In the context of AI for medical imaging, what does the ``Clever Hans'' effect
refer to?
\begin{itemize}[leftmargin=1.8em]
  \choice A bias that arises when up-scaling the resolution of low-quality medical images.
  \choice A case where a model looks like it performs well but is actually leaning on irrelevant features or artifacts in the data.
  \choice A case where an AI predicts outcomes incorrectly because its training data were biased.
  \choice The use of data augmentation to make medical-imaging models more robust.
\end{itemize}
\closequestion

\questiontitle{Explainable AI and interpretability in medical imaging}
Which statements about explainable AI and interpretability in medical imaging are
correct?
\begin{itemize}[leftmargin=1.8em]
  \choice Gradient integration propagates a relevance value backwards through the layers to the input, estimating each input feature's relevance.
  \choice Explainable-AI methods such as Grad-CAM help visualise which parts of a medical image matter most for a prediction.
  \choice Occlusion analysis deletes patches from the training images to make the model more interpretable.
  \choice Counterfactual explanations generate alternative images showing what would have had to change for a different diagnosis.
\end{itemize}
\closequestion

\questiontitle{AlphaFold 2}
Which of the following statements about AlphaFold 2 are correct?
\begin{itemize}[leftmargin=1.8em]
  \choice AlphaFold 2 relies on quantum-mechanical calculations to sharpen its predictions.
  \choice AlphaFold 2 draws on multiple sequence alignments and homologous protein structures to improve its accuracy.
  \choice AlphaFold 2 is a deep-learning model that predicts a protein's three-dimensional structure from its amino-acid sequence.
  \choice AlphaFold 2 is built on a convolutional neural network architecture.
\end{itemize}
\closequestion


% ==========================================================
%  SECTION 6 — Biological Sequences
% ==========================================================
\section{Biological Sequences}

\questiontitle{Composition of biological sequences}
Which statements about the composition of biological sequences are correct?
\begin{itemize}[leftmargin=1.8em]
  \choice Protein sequences are chains of amino acids joined by peptide bonds.
  \choice DNA sequences are built from nucleic acids.
  \choice Biological sequences obey no particular patterns or rules in their make-up.
  \choice RNA sequences are built from lipids.
\end{itemize}
\closequestion

\questiontitle{Deep learning for biological sequences}
Which of the following statements about deep learning methods for biological
sequences are correct?
\begin{itemize}[leftmargin=1.8em]
  \choice Feature extraction from biological sequences can reveal motifs, i.e.\ recurring patterns or elements within them.
  \choice One-hot encoding is a common way to prepare biological sequences as input for deep models.
  \choice Deep-learning techniques apply to tasks such as classification, segmentation, and generative modelling of biological sequences.
  \choice 1D convolutions, Recurrent Neural Networks (RNNs), Long Short-Term Memory networks (LSTMs), and Transformers are all used to analyse biological sequences.
\end{itemize}
\closequestion


% ==========================================================
%  2026 EXAM (single-best-answer) — Sections 7–10
% ==========================================================
\begin{notebox}[2026 exam --- single-best-answer]
The remaining sections are from the \textbf{2026} exam (25 June 2026). Each of these
questions has \textbf{exactly one} correct option among the three listed. Both the
question and the options have been paraphrased; the meaning and correct answers are
preserved.
\end{notebox}


% ==========================================================
%  SECTION 7 — Molecules, Representations & Cheminformatics (2026)
% ==========================================================
\section{Molecules, Representations \& Cheminformatics (2026)}

\questiontitle{Molecular graphs vs.\ fingerprints}
Which of the following statements accurately describes the relationship between
molecular graphs and molecular fingerprints?
\begin{itemize}[leftmargin=1.8em]
  \choice A fingerprint can be computed from a molecular graph, yet the exact graph usually cannot be recovered from the fingerprint.
  \choice Fingerprints and molecular graphs must hold the same structural information.
  \choice A molecular graph can always be reconstructed exactly from its fingerprint.
\end{itemize}
\closequestion

\questiontitle{Converting between SMILES and fingerprints}
Using standard cheminformatics tools, in which direction is conversion between a
SMILES string and a molecular fingerprint actually possible?
\begin{itemize}[leftmargin=1.8em]
  \choice Cheminformatics tools can convert SMILES and fingerprints into one another in both directions.
  \choice Cheminformatics tools can turn a SMILES into a fingerprint, but not the other way round.
  \choice Cheminformatics tools can turn a fingerprint into a SMILES, but not the other way round.
\end{itemize}
\closequestion

\questiontitle{Tanimoto coefficient}
On what quantity is the Tanimoto coefficient based when it scores how similar two
molecular fingerprints are?
\begin{itemize}[leftmargin=1.8em]
  \choice It takes the count of shared features over the total number of features present in either molecule.
  \choice It superimposes three-dimensional pharmacophore points and averages their spatial distances.
  \choice It compares the core scaffolds, scoring one whenever the scaffolds coincide.
\end{itemize}
\closequestion

\questiontitle{Aspirin SMILES}
Which of the following SMILES strings represents aspirin (acetylsalicylic acid)?
\begin{itemize}[leftmargin=1.8em]
  \choice \texttt{O=C(C)Oc1ccccc1C(=O)O}
  \choice \texttt{cc(=O)Oc1ccccc1c(=O)O}
  \choice \texttt{O=C[C]Oc1ccccc1C[=O]O}
\end{itemize}
\closequestion

\questiontitle{Bioactivity datasets}
Which of the following best captures what bioactivity datasets usually look like in practice?
\begin{itemize}[leftmargin=1.8em]
  \choice They are usually large, balanced, and almost free of measurement noise.
  \choice They are frequently small, imbalanced, and noisy, with far more inactive than active compounds.
  \choice They are cheap to produce and therefore cover most compound--assay combinations.
\end{itemize}
\closequestion

\questiontitle{The Design--Make--Test--Analyse cycle}
In drug discovery, which option correctly captures how the Design--Make--Test--Analyse
(DMTA) cycle proceeds?
\begin{itemize}[leftmargin=1.8em]
  \choice Compounds are tested before being designed, and only the active ones are afterwards synthesized.
  \choice Compounds are designed and made just once, after which testing removes any need for further design.
  \choice Compounds are designed, made, and tested experimentally, and the analysis steers the next design round.
\end{itemize}
\closequestion


% ==========================================================
%  SECTION 8 — Virtual Screening, QSAR & Targets (2026)
% ==========================================================
\section{Virtual Screening, QSAR \& Targets (2026)}

\questiontitle{Structure-based vs.\ ligand-based drug design}
What is the essential difference between structure-based drug design (SBDD) and
ligand-based drug design (LBDD)?
\begin{itemize}[leftmargin=1.8em]
  \choice SBDD needs known active ligands, whereas LBDD needs the target's 3D structure.
  \choice Both SBDD and LBDD need a high-resolution 3D structure of the target.
  \choice SBDD relies on the target's 3D structure, whereas LBDD can work from known ligands without that structure.
\end{itemize}
\closequestion

\questiontitle{Similarity search vs.\ model-based search}
In what way do similarity search and model-based search differ from each other?
\begin{itemize}[leftmargin=1.8em]
  \choice Similarity search is ligand-based, whereas model-based search must have the 3D structure of the target's binding site.
  \choice Similarity search retrieves molecules resembling a query, whereas model-based search predicts activity with a model trained on labelled data.
  \choice Similarity search learns a decision function from labelled actives and inactives, whereas model-based search ranks compounds by closeness to a query.
\end{itemize}
\closequestion

\questiontitle{Early-recognition metrics (odd one out)}
Of the metrics below, which one is \emph{not} built to reward finding actives early,
i.e.\ near the top of a ranked virtual-screening list?
\begin{itemize}[leftmargin=1.8em]
  \choice BEDROC.
  \choice AUC-ROC.
  \choice Enrichment factor.
\end{itemize}
\closequestion

\questiontitle{Evaluating a virtual-screening model}
A virtual-screening model orders a million compounds, yet your budget only permits
assaying the top 1\%. Which model-selection strategy best fits this situation?
\begin{itemize}[leftmargin=1.8em]
  \choice Pick the model with the best accuracy at a 0.5 probability threshold, since ranking above that threshold no longer matters once compounds are classified.
  \choice Pick the model with the highest overall AUC-ROC, since weighting all ranking positions equally guarantees the best enrichment in the top 1\%.
  \choice Favour an early-recognition metric such as the enrichment factor near the 1\% cutoff or BEDROC, since the top of the ranking decides experimental usefulness.
\end{itemize}
\closequestion

\questiontitle{Defensible generalization estimates in QSAR}
In a QSAR pipeline, which procedure yields the most trustworthy estimate of how well
the model generalises?
\begin{itemize}[leftmargin=1.8em]
  \choice Fix the split first, carry out feature and hyperparameter selection on the training data, and reserve the held-out test set for the final evaluation only.
  \choice Split first, then compare models and hyperparameters on the test set and report the best test-set score.
  \choice Select features on the whole dataset, then split the reduced data and tune hyperparameters on the training part only.
\end{itemize}
\closequestion

\questiontitle{Compound-series bias}
What does compound-series bias refer to, and why does it undermine QSAR model evaluation?
\begin{itemize}[leftmargin=1.8em]
  \choice Because many compounds are closely related, a random split can put near-analogues in both training and test, underestimating the generalization error.
  \choice Test-compound information leaks into feature or hyperparameter selection, biasing the reported score even when compounds are structurally unrelated.
  \choice The active class is far smaller than the inactive one, so the majority class dominates performance even when the chemical series are diverse.
\end{itemize}
\closequestion

\questiontitle{Drug synergy and model symmetry}
Which option combines a correct definition of drug synergy with a sensible property
that a model predicting it should have?
\begin{itemize}[leftmargin=1.8em]
  \choice The combined effect is larger than the additive one, and the prediction should stay the same when the two compounds are swapped.
  \choice The combined effect equals the additive expectation, and the prediction may differ when the pair is supplied in reverse order.
  \choice The combined effect is below the additive expectation, and the prediction should be unaffected by the order of the compounds.
\end{itemize}
\closequestion

\questiontitle{Proteochemometric modelling}
In what respect does proteochemometric modelling go beyond classical single-target QSAR?
\begin{itemize}[leftmargin=1.8em]
  \choice It uses representations of both the molecule and the protein target together.
  \choice It uses molecular representations with a softmax over the candidate protein targets.
  \choice It uses molecular representations with one separate prediction head per protein target.
\end{itemize}
\closequestion

\questiontitle{Target deconvolution / target fishing}
How is the task known as target deconvolution (or target fishing) defined?
\begin{itemize}[leftmargin=1.8em]
  \choice Given one protein and a compound library, rank every molecule by predicted affinity to that protein.
  \choice Given one molecule and several candidate proteins, predict which protein it is most likely to bind.
  \choice Given one molecule, predict binding to each candidate protein independently, so any number of targets may be assigned at once.
\end{itemize}
\closequestion


% ==========================================================
%  SECTION 9 — Graph Neural Networks (2026)
% ==========================================================
\section{Graph Neural Networks (2026)}

\questiontitle{Reordering atoms in a molecular graph}
Suppose the atoms of a molecular graph are given new indices while the bonds between
them are left unchanged. What behaviour do we expect from a graph neural network?
\begin{itemize}[leftmargin=1.8em]
  \choice The node representations are permuted accordingly, but the graph-level prediction may change because the adjacency matrix was reordered.
  \choice Both the node representations and the graph-level prediction stay in exactly the same tensor positions.
  \choice The node representations are permuted accordingly, while the graph-level prediction stays unchanged.
\end{itemize}
\closequestion

\questiontitle{Steps of an MPNN layer}
What is the usual order of operations carried out inside a single Message Passing
Neural Network (MPNN) layer?
\begin{itemize}[leftmargin=1.8em]
  \choice Generate messages, aggregate them, then update.
  \choice Generate messages, apply graph convolution, then aggregate.
  \choice Apply graph convolution, generate messages, then update.
\end{itemize}
\closequestion

\questiontitle{Oversmoothing in GNNs}
Once a large number of message-passing layers are stacked, the node embeddings across
a molecular graph end up nearly identical. Why does this cause trouble?
\begin{itemize}[leftmargin=1.8em]
  \choice Long-range information is squeezed through narrow graph bottlenecks, so distant signals cannot travel effectively.
  \choice The model may no longer tell local atom environments apart because the node embeddings grow too similar.
  \choice Important information never reaches distant nodes because too few message-passing steps were taken.
\end{itemize}
\closequestion

\questiontitle{Contrastive learning}
By what mechanism does contrastive learning build a joint molecule--target embedding
space?
\begin{itemize}[leftmargin=1.8em]
  \choice It clusters molecules and targets separately and assumes equal-index clusters are binding pairs.
  \choice It pulls associated molecule-target pairs closer together and pushes mismatched pairs apart.
  \choice It concatenates molecule and target representations and fits a supervised activity regressor without comparing pairs in a shared space.
\end{itemize}
\closequestion

\questiontitle{Multiple Instance Learning}
How does Multiple Instance Learning depart from ordinary supervised learning, in
which every feature vector carries a single label $(\boldsymbol{x}, y)$?
\begin{itemize}[leftmargin=1.8em]
  \choice \ldots we are given a bag of feature vectors that share a single label $(\{\boldsymbol{x}_1, \ldots, \boldsymbol{x}_N\}, y)$.
  \choice \ldots we are given a bag of feature vectors paired with a bag of labels $(\{\boldsymbol{x}_1, \ldots, \boldsymbol{x}_N\}, \{y_1, \ldots, y_M\})$.
  \choice \ldots we are given single feature vectors paired with a bag of labels $(\boldsymbol{x}, \{y_1, \ldots, y_M\})$.
\end{itemize}
\closequestion


% ==========================================================
%  SECTION 10 — Generative Models & Synthesis Planning (2026)
% ==========================================================
\section{Generative Models \& Synthesis Planning (2026)}

\questiontitle{Reward hacking a QSAR proxy}
A generator is trained to maximise a QSAR-based activity score and learns to output
molecules the score rates as highly active, yet later wet-lab assays find them
inactive. What most plausibly went wrong?
\begin{itemize}[leftmargin=1.8em]
  \choice The activity score was supplied only after each molecule was finished, so the reward was delayed.
  \choice The generator gamed weaknesses in the proxy scoring model instead of raising genuine experimental activity.
  \choice The generator made too few structurally distinct molecules, lowering the diversity of candidates.
\end{itemize}
\closequestion

\questiontitle{Combating analogue rediscovery (odd one out)}
A goal-directed generator keeps producing close analogues of molecules it has already
scored highly. Which of the following would \emph{not} help to counteract this
tendency?
\begin{itemize}[leftmargin=1.8em]
  \choice Blend samples from a fixed exploratory generator with those of the trainable generator while generating SMILES.
  \choice Penalise the reward of molecules resembling earlier high-scoring ones kept in a memory buffer.
  \choice Push exploitation harder by always taking the highest-scoring molecule as the next training example.
\end{itemize}
\closequestion

\questiontitle{Autoregressive molecular reinforcement learning}
When molecule construction is performed autoregressively and treated as a
reinforcement-learning problem, which assignment of the RL roles is correct?
\begin{itemize}[leftmargin=1.8em]
  \choice The partial molecule is the state and the next building step is the action, but the scoring function is the policy and the generator gives the reward.
  \choice The finished molecule is the state, the scoring function is the action, and each partial molecule gets its own final reward.
  \choice The partial molecule is the state, the next building step is the action, the generator is the policy, and the finished molecule earns the reward.
\end{itemize}
\closequestion

\questiontitle{Diffusion vs.\ one-shot latent decoding}
What sets diffusion-based molecular generation apart from decoding a molecule from a
single latent vector in one pass?
\begin{itemize}[leftmargin=1.8em]
  \choice It maps a single latent vector straight to a finished molecule in one decoder pass.
  \choice It begins from noise and repeatedly applies a learned denoising process.
  \choice It builds the molecule step by step, predicting the next token or graph action from the partial molecule.
\end{itemize}
\closequestion

\questiontitle{Conditional generation}
Trained on molecule--logP pairs, a model takes a target logP value as input and emits
molecules in a single forward pass, never querying a logP scoring function during
generation. Which generative paradigm is this?
\begin{itemize}[leftmargin=1.8em]
  \choice Conditional generation.
  \choice Goal-directed generation.
  \choice Distribution learning followed by unconditional sampling.
\end{itemize}
\closequestion

\questiontitle{Distribution learning}
In molecular generation, which option most accurately captures what distribution learning is?
\begin{itemize}[leftmargin=1.8em]
  \choice Distribution learning produces molecules conditioned on side information such as protein targets or desired property values.
  \choice Distribution learning seeks to produce molecules resembling the training set by learning its underlying distribution.
  \choice Distribution learning concentrates on producing molecules optimised for specific biological or physicochemical properties.
\end{itemize}
\closequestion

\questiontitle{Mode collapse}
In a generative model for molecules, what does the term mode collapse refer to?
\begin{itemize}[leftmargin=1.8em]
  \choice The model produces molecules that are valid and novel but match the training-set property distribution poorly.
  \choice The model keeps emitting the same or nearly the same molecules.
  \choice The model emits almost only molecules that already appear in the training set.
\end{itemize}
\closequestion

\questiontitle{Reaction prediction and retrosynthesis}
Which of the options correctly pairs each task with its corresponding inputs and outputs?
\begin{itemize}[leftmargin=1.8em]
  \choice Both forward prediction and single-step retrosynthesis take only the reaction conditions and predict the yield.
  \choice Forward prediction maps reactants to products, while single-step retrosynthesis maps a target product to candidate precursors.
  \choice Forward prediction maps products to reactants, while single-step retrosynthesis predicts the major product from given reactants.
\end{itemize}
\closequestion

\questiontitle{Reaction templates}
What role is a reaction template intended to serve?
\begin{itemize}[leftmargin=1.8em]
  \choice To give a general symbolic rule stating the required substructures and how reactants turn into products.
  \choice To capture reaction conditions such as solvent and temperature, leaving the structural change to a neural model.
  \choice To attach a reaction-class label grouping similar reactions, without saying how atoms or bonds change.
\end{itemize}
\closequestion

\questiontitle{Easiest reaction task}
Among the tasks listed below, which one is typically the least difficult?
\begin{itemize}[leftmargin=1.8em]
  \choice Single-step retrosynthesis.
  \choice Multi-step retrosynthesis planning.
  \choice Forward reaction prediction.
\end{itemize}
\closequestion

\questiontitle{Template-free translation methods}
In what terms do template-free translation methods formulate reaction prediction or
retrosynthesis?
\begin{itemize}[leftmargin=1.8em]
  \choice As looking up the most compatible stored reaction pattern in an associative memory.
  \choice As sorting the input into a learned reaction category and then applying a symbolic transformation.
  \choice As a direct sequence-to-sequence translation between molecular reaction strings.
\end{itemize}
\closequestion

\questiontitle{Deep networks inside MCTS for retrosynthesis}
What makes deep neural networks helpful when they are embedded within Monte Carlo
Tree Search for retrosynthesis?
\begin{itemize}[leftmargin=1.8em]
  \choice They take over the search by scoring only the first disconnection and then running that route greedily.
  \choice They can estimate the value of unvisited states and propose promising actions when exhaustive search is impractical.
  \choice They mostly balance chemical equations at each node, while the search alone fixes value and policy without learned guidance.
\end{itemize}
\closequestion

\questiontitle{Pocket-conditioned 3D ligand generation}
Which option most accurately describes pocket-conditioned 3D ligand generation?
\begin{itemize}[leftmargin=1.8em]
  \choice Ligands are built without protein context and only docked afterwards; the docking score does not steer the generator.
  \choice The protein pocket supplies geometric context for building or denoising a ligand, usually followed by validity, docking, property, and synthesis checks.
  \choice A fixed ligand graph is supplied and only different conformations of that same molecule are sampled, ignoring the protein.
\end{itemize}
\closequestion


% ==========================================================
%  ANSWER KEY
% ==========================================================
\newpage
\section*{Answer Key}
\addcontentsline{toc}{section}{Answer Key}

\begin{keybox}[Note on scoring]
Questions 1--34 (the 2024 exam) are multiple-select: each listed letter is a
correct option, and any option \emph{not} listed is incorrect. On the original
Moodle exam, each correct tick added points and each incorrect tick subtracted
them. Questions 35--67 (the 2026 exam) are single-best-answer: exactly one letter
is correct.
\end{keybox}

\begin{enumerate}[label=\textbf{Q\arabic*.}, leftmargin=3em, itemsep=4pt]
  % --- Section 1 ---
  \item \textbf{A, B} --- SMILES is the \emph{Simplified Molecular Input Line Entry System} and is a linear character sequence. It \emph{can} be converted to a graph, and is \emph{not} unique (one molecule has many valid SMILES).
  \item \textbf{A, B} --- both \texttt{O=C(C(=O)O)C} and \texttt{CC(=O)C(=O)O} encode CH\textsubscript{3}--CO--COOH. \texttt{CCCOOO} and \texttt{OCCOO} do not.
  \item \textbf{B, C} --- fingerprints are binary/numerical descriptors used for similarity search and ML. They are lossy (no unique graph reconstruction) and carry no 3D information.
  \item \textbf{A, B, C} --- bioassays measure a biological response, test at various concentrations (dose--response), and are wet-lab procedures. ``\emph{in vitro} = in a living organism'' is false (\emph{in vivo} is the living organism).
  \item \textbf{A, C, D} --- Lipinski's Ro5 is an ADME proxy, clearance is hard to model, and ADME is central to pharmacokinetics. Good ADME does \emph{not} guarantee an effective drug.
  \item \textbf{C, D} --- QSAR data often comes in structurally related series and is noisy. It is usually \emph{small} (not thousands/millions) and \emph{imbalanced} (not balanced).
  \item \textbf{B, C} --- QSAR predicts drug candidates and toxicity. 3D protein modeling and genome mapping are different problems.
  % --- Section 2 ---
  \item \textbf{A} --- fingerprints are standard in ligand-based VS. It \emph{is} suitable for unknown compounds, does \emph{not} need the protein structure, and considers more than bioactivity (ADME, filters).
  \item \textbf{A, D} --- temporal splits simulate real-world acquisition; cluster splits keep structurally similar compounds out of both sets. Random splits do \emph{not} eliminate bias, and CV \emph{is} used.
  \item \textbf{B, D} --- BEDROC and the Enrichment Factor reward early recognition. Precision-Recall curve and Average Precision are general metrics, not early-recognition-specific.
  \item \textbf{A, B} --- target identification finds the \emph{most likely} target via multi-class classification. ``All proteins'' and ``multi-task'' describe target \emph{prediction}, not identification.
  \item \textbf{A, B} --- synergy is ``greater than the sum''; the Loewe additivity model assesses it. Antagonism is \emph{less} than the sum, and synergy applies to any drug pair (not only similar ones).
  \item \textbf{B, C} --- RF overfits less and is more interpretable than GNNs. GNNs are \emph{not} universally better and are \emph{not} less sensitive to data quality.
  % --- Section 3 ---
  \item \textbf{A, B} --- a molecular graph is a node-feature matrix plus an edge/adjacency matrix; node features include atomic number, hybridization, charge. GNNs \emph{do} use edge information, and bond type \emph{can} be encoded.
  \item \textbf{B, C, D} --- GNNs struggle with long-range interactions, very large graphs, and are compute-heavy. They are \emph{not} ``incapable of predicting any property accurately''.
  \item \textbf{A, B} --- multi-modal contrastive learning maps modalities into a shared latent space and is \emph{self-supervised}. It does not ``concentrate representations to make predictions'', and is \emph{not} unsupervised.
  \item \textbf{C, D} --- a MIL bag has a \emph{single} label and the model must be permutation-invariant. Instances are \emph{not} individually labelled, and bag size need \emph{not} be fixed.
  % --- Section 4 ---
  \item \textbf{D} --- autoregressive models can generate SMILES. Modern generators produce mostly valid SMILES, generative models explore \emph{novel} space, and GANs \emph{can} generate molecules.
  \item \textbf{C, D} --- goal-directed generation targets \emph{specific} desired properties and can be optimised with RL. Replicating existing structures / matching the training distribution describes \emph{distribution learning}.
  \item \textbf{B, D} --- FCD compares means and covariances of \emph{latent} (ChemNet) representations and measures distance between \emph{distributions} of generated vs.\ real molecules --- not individual atoms, not raw structures.
  \item \textbf{A, B, C} --- reaction templates, training on known-reaction databases, and transformer architectures are all used. QM data is \emph{not} required.
  \item \textbf{A, B} --- synthesis planning designs a reaction sequence and uses retrosynthetic analysis. It is \emph{not} trivial, and involves more than just picking reactants.
  \item \textbf{A, C, D} --- per the lecture's framing the \emph{action} is selecting reactants, the \emph{environment} applies the reaction and returns new molecules, and the \emph{reward} comes when the current set are all available building blocks. Only the \emph{state} option is wrong: the state is the \emph{current set of molecules} to be made, not the set of available building blocks.
  % --- Section 5 ---
  \item \textbf{A} --- cholesterol is numeric, so the \emph{median} handles skewed missing values. Smoking status is categorical (not continuous/mean-imputable), and cholesterol is continuous, not merely ordinal.
  \item \textbf{A, C} --- ignored batch effects make models separate batches instead of biology, and they create artificial differences. They are \emph{not} removed by more samples and are \emph{not} beneficial.
  \item \textbf{C, D} --- changing prevalence is a label/prior shift $p(y)$; changing the diagnostic procedure is a concept shift $p(y\mid x)$. Good validation accuracy does \emph{not} auto-adapt, and $p(x)$ \emph{can} shift.
  \item \textbf{B, C} --- X-rays distinguish tissue densities; MRI uses magnetic fields and radio waves. CT uses X-rays (not ultrasound), and OCT uses light (not X-rays).
  \item \textbf{B, C} --- classification and segmentation are the dominant medical-imaging tasks (segmentation especially, vs.\ typical vision tasks).
  \item \textbf{B, D} --- U-Net has a symmetric encoder/decoder with skip connections (encoder feature maps concatenated into the decoder). It is convolutional (not attention-based) and was built for \emph{segmentation}, not detection.
  \item \textbf{B} --- the Clever Hans effect is a model that \emph{looks} good but relies on irrelevant features/artifacts (shortcut learning).
  \item \textbf{B, D} --- Grad-CAM highlights important image regions; counterfactuals show what change would flip the diagnosis. The first option mislabels LRP as ``gradient integration'', and occlusion perturbs \emph{input}, not training data.
  \item \textbf{B, C} --- AlphaFold~2 uses multiple sequence alignments + homologous structures to predict 3D structure from sequence. It uses an attention-based (Evoformer) architecture, \emph{not} a CNN, and no QM calculations.
  % --- Section 6 ---
  \item \textbf{A, B} --- proteins = amino acids via peptide bonds; DNA = nucleic acids. Sequences \emph{do} follow rules, and RNA is nucleic acid, not lipid.
  \item \textbf{A, B, C, D} --- all true: motif/feature extraction, one-hot input encoding, classification/segmentation/generation tasks, and 1D-CNN / RNN / LSTM / Transformer architectures.
  % --- Section 7 (2026) ---
  \item \textbf{A} --- a fingerprint is computed from the graph but is lossy/hashed, so the graph cannot be recovered; ``same information'' (B) and ``exact reconstruction'' (C) are false.
  \item \textbf{B} --- SMILES $\to$ fingerprint is standard; fingerprints are lossy, so fingerprint $\to$ SMILES is impossible.
  \item \textbf{A} --- Tanimoto $=$ shared bits divided by the union of bits (intersection over union). B is 3D pharmacophore alignment; C is scaffold matching.
  \item \textbf{A} --- \texttt{O=C(C)Oc1ccccc1C(=O)O} is valid acetylsalicylic acid; B misuses lowercase (aromatic) atoms for the carbonyls; C misuses square brackets.
  \item \textbf{B} --- bioactivity data is limited, imbalanced, and noisy, with far more inactives than actives.
  \item \textbf{C} --- design, make (synthesize), test, analyse, then iterate to the next design.
  % --- Section 8 (2026) ---
  \item \textbf{C} --- SBDD uses the target's 3D structure; LBDD uses known ligands without it.
  \item \textbf{B} --- similarity search retrieves neighbours of a query; model-based search uses a supervised model learned from labels.
  \item \textbf{B} --- AUC-ROC weights the whole ranking equally; BEDROC and the enrichment factor are early-recognition metrics.
  \item \textbf{C} --- use an early-recognition metric (enrichment factor near the 1\% cutoff, or BEDROC).
  \item \textbf{A} --- split first, select features/hyperparameters on training data only, then evaluate on the held-out test set exactly once.
  \item \textbf{A} --- structurally related analogues split across train and test under a random split, giving an over-optimistic estimate.
  \item \textbf{A} --- synergy is super-additive, and a synergy-prediction model should be order-invariant (symmetric in the two compounds).
  \item \textbf{A} --- it jointly represents molecules \emph{and} protein targets (single-target QSAR uses the molecule alone).
  \item \textbf{B} --- given a molecule and several candidate proteins, predict the most likely target.
  % --- Section 9 (2026) ---
  \item \textbf{C} --- node representations permute with the relabelling (equivariance), while the graph-level prediction is unchanged (invariance).
  \item \textbf{A} --- message generation, aggregation, update.
  \item \textbf{B} --- oversmoothing: node embeddings become too similar, so local atom-environment information is lost (A is over-squashing; C is too few layers).
  \item \textbf{B} --- pull associated molecule--target pairs together and push mismatched pairs apart.
  \item \textbf{A} --- a bag of feature vectors carrying a single shared label.
  % --- Section 10 (2026) ---
  \item \textbf{B} --- reward hacking: the generator exploited weaknesses of the proxy score rather than improving true activity.
  \item \textbf{C} --- always picking the single highest-scoring molecule increases exploitation and worsens rediscovery, so it is \emph{not} a remedy (A and B both add diversity).
  \item \textbf{C} --- partial molecule $=$ state, next step $=$ action, generator $=$ policy, completed molecule $=$ reward.
  \item \textbf{B} --- diffusion starts from noise and iteratively denoises (A is one-shot decoding; C is autoregressive).
  \item \textbf{A} --- conditional generation: conditioned on the target property, one pass, no scoring function queried.
  \item \textbf{B} --- distribution learning models the training distribution and generates molecules resembling it.
  \item \textbf{B} --- mode collapse: the model repeatedly produces the same or very similar molecules.
  \item \textbf{B} --- forward prediction maps reactants $\to$ products; single-step retrosynthesis maps product $\to$ precursors.
  \item \textbf{A} --- a reaction template is a generalized symbolic rule giving required substructures and the reactant $\to$ product transformation.
  \item \textbf{C} --- forward reaction prediction is generally the easiest (single-step retro is one-to-many; multi-step planning is hardest).
  \item \textbf{C} --- template-free methods frame the task as direct sequence-to-sequence translation between reaction strings.
  \item \textbf{B} --- deep networks estimate the value of unvisited states and suggest promising actions when exhaustive search is infeasible.
  \item \textbf{B} --- the binding pocket provides geometric context to generate/denoise the ligand, usually followed by validity/docking/property/synthesis checks.
\end{enumerate}


\end{document}
